\documentclass[conference]{IEEEtran}

\usepackage{amsmath,amssymb}
\usepackage{booktabs}
\usepackage{cite}
\usepackage{graphicx}
\usepackage[hidelinks]{hyperref}

\title{Microchip Character Recognition Using Image Processing and Machine Learning}

\author{
  \IEEEauthorblockN{Group 8}
  \IEEEauthorblockA{
    ME5411 Robot Vision and AI\\
    National University of Singapore\\
    Singapore
  }
}

\begin{document}

\maketitle

\begin{abstract}
This report presents an image-processing and machine-learning pipeline for
recognizing characters printed on a microchip. The final report will describe
the preprocessing, segmentation, feature extraction, classification, and
evaluation procedures developed by Group 8.
\end{abstract}

\begin{IEEEkeywords}
robot vision, image processing, character recognition, convolutional neural
network, classification
\end{IEEEkeywords}

\section{Introduction}
Character recognition in low-contrast industrial imagery requires careful
preprocessing, segmentation, and classification. Classical image-processing
methods provide interpretable tools for enhancing and isolating characters
\cite{gonzalez2018digital}, while convolutional neural networks can learn
features directly from image data \cite{lecun1998gradient}.

The implementation is organised as six reproducible tasks. Tasks 1--4 use
classical image-processing operators covered in the lectures; Tasks 5 and 6
compare a compact CNN with a conventional HOG-based classifier. Since the
course images are kept outside the repository, numerical results are generated
by the supplied commands after the data are placed in \texttt{data/raw/}.

\section{Problem Formulation}
Describe the input image, recognition targets, operating assumptions, and
evaluation criteria. State the dataset split and measures used for comparison.

\section{Methodology}

\subsection{Image Enhancement and Smoothing}
Task 1 first converts each input to an 8-bit grayscale image and applies the
piecewise-linear contrast transform
\begin{equation}
  g(x,y)=\operatorname{clip}\left(255\frac{f(x,y)-r_{\min}}
  {r_{\max}-r_{\min}},0,255\right).
  \label{eq:stretch}
\end{equation}
The input limits are measured from the image unless explicitly supplied. A
histogram-equalisation alternative is also implemented for images whose
histogram occupies only a narrow range. Task 2 provides the normalized
$3\times3$ (or arbitrary odd-size) averaging mask from the lectures, a
separable Gaussian mask, a median filter for impulse noise, and the rotating
mask average that chooses the locally most homogeneous window. The selected
filter and size are recorded in the command line used for an experiment.

\subsection{Region of Interest and Binarization}
Task 3 accepts an explicit $(x,y,w,h)$ text-line ROI so that camera-specific
geometry is not hidden in the algorithm. When no ROI is supplied the full image
is used. Otsu's between-class variance selects a global threshold automatically;
local mean thresholding is available when illumination varies across the chip.
Both bright-on-dark and dark-on-bright text are supported. Foreground pixels are
stored as one (and written as 255 in PNG files).

\subsection{Outline Extraction and Character Segmentation}
Task 4 computes the horizontal and vertical Sobel responses and their gradient
magnitude, yielding the outline diagnostic required by the segmentation
lectures. The binary foreground is then labelled with either 4- or
8-connectivity using a flood-fill implementation. Components smaller than the
configurable area threshold are rejected as noise. Each remaining component is
sorted left-to-right, padded, and exported as a character crop together with
its bounding box, centroid, area, and perimeter in JSON. These diagnostics make
merged or split characters visible before model training.

\subsection{CNN Classifier}
Task 5 resizes each crop to $32\times32$ and scales intensities to $[0,1]$.
The network has three convolutional blocks (16, 32, and 64 channels), batch
normalisation/ReLU, max pooling, adaptive average pooling, and two fully
connected layers. Cross-entropy loss is optimized with Adam (learning rate
$10^{-3}$ by default), using a fixed seed and a stratified 80/20 split. A
checkpoint stores the class names, preprocessing size, seed, and epoch history.

\subsection{Non-CNN Classifier}
Task 6 uses a 32-by-32 unsigned-gradient HOG descriptor (4-pixel cells and nine
orientation bins). An RBF SVM with $C=10$ is used when scikit-learn is present;
the same interface falls back to a nearest-centroid classifier for a minimal
installation. The feature extractor, split seed, classifier type, and accuracy
are serialized so that the baseline can be compared with the CNN without
changing the data protocol.

\section{Experimental Setup}
Report the software environment, hardware, random seeds, dataset split, data
preprocessing, hyperparameter search, and evaluation metrics. Ensure that all
experiments can be reproduced using the submitted code.

\section{Results}
Present quantitative and qualitative results. Include intermediate
image-processing outputs, classification metrics, confusion matrices, runtime
measurements, and sensitivity studies.

\begin{table}[t]
  \caption{Placeholder Classification Results}
  \label{tab:results}
  \centering
  \begin{tabular}{lcc}
    \toprule
    Method & Accuracy (\%) & Runtime (ms) \\
    \midrule
    CNN & -- & -- \\
    Non-CNN & -- & -- \\
    \bottomrule
  \end{tabular}
\end{table}

\section{Discussion}
Compare the effectiveness, efficiency, robustness, and failure modes of the two
classification approaches. Explain observed differences using evidence from the
experiments rather than accuracy alone.

\section{Conclusion}
Summarize the principal findings, lessons learned, limitations, and possible
improvements.

\section*{Code Execution Guide}
Install the dependencies listed in \texttt{requirements.txt}, then run the
following from the repository root with \texttt{PYTHONPATH=src}:

\begingroup
\scriptsize
\begin{verbatim}
python -m task_01.enhance --input data/raw/chip.png \
  --output results/figures/enhanced.png
python -m task_02.smooth --input results/figures/enhanced.png \
  --output results/figures/smoothed.png --method gaussian --size 3
python -m task_03.threshold --input results/figures/smoothed.png \
  --output results/figures/binary.png --method otsu
python -m task_04.segment --input results/figures/binary.png \
  --output-dir data/processed/characters
python -m task_05.cnn --data data/processed/characters \
  --model results/models/character_cnn.pt
python -m task_06.classical train --data data/processed/characters \
  --model results/models/hog_svm.pkl
\end{verbatim}
\endgroup

The classifier data set has one directory per character (for example
\texttt{characters/A/*.png}). Task 4 crops are assigned to these directories
using the available ground truth before training. Every command accepts a
\texttt{--help} page with parameter defaults and writes generated artifacts only
under \texttt{results/} or \texttt{data/processed/}.

\section*{Acknowledgment}
List external libraries, toolboxes, third-party code, and other assistance that
must be acknowledged according to the course requirements.

\bibliographystyle{IEEEtran}
\bibliography{references}

\end{document}
